% translation-source: docs/user/latex/chapters/16-depannage.tex; sha256: cbf3419c9b5a3f8fab807eacd480befd6b0240f42f1f53c0d66a95c7a6dc7021
\chapter{Checks and troubleshooting}
\section{A command seems unavailable}
Check that a photo is open, wait for processing and confirm it is a RAW for
camera-development commands. A missing profile is not necessarily a list error:
availability depends on installed files and compatibility.

\section{The photo differs from another application}
Compare the same file, framing and neutral state first. Check camera profile,
white balance, calibration, optical corrections, noise and sharpening.
An identical displayed temperature does not alone establish equal development.
The viewing background and color-management path may also affect comparison.

\section{Settings reappear on opening}
Look for the \chemin{.ipr} sidecar kept with the photo. Restoration is intentional.
Use the reset appropriate to your goal instead of immediately deleting files.
Keep a backup before manually changing a sidecar or configuration.

\section{Saving or export fails}
Read the bottom status message. Check free disk space, folder permissions,
availability and file locks. A sidecar failure is not a successful export.
The export success message gives the actual path, including any added number.

\section{A function is not described}
Crop, masks and AI panels are reserved in this edition. Complete filmstrip
badges, comprehensive modified-group indicators and detailed engine states
remain future work, not hidden commands. User visual validation and testing
other cameras remain necessary despite automated interface tests.
